% =========================================================================
% §2 Related Work
% Budget: ~1.0 pp
% Drafting order: After Discussion (so we know what we positioned against).
%
% Content blueprint:
%   - Coding-agent harnesses: SWE-Agent, OpenHands/OpenDevin, Aider,
%     AutoCodeRover, OpenCode, SuperCoder. One line each.
%   - Retrieval approaches: agentic grep (Claude Code), file-dependency
%     PageRank (Aider), semantic search (LocAgent, RepoGraph, CGM),
%     persistent structural index (this work + Augment / code-review-graph).
%   - Localization metrics: LocAgent / Agentless / SWE-Explore / TRAJEVAL.
%     Note the field's move toward stage-decomposed trajectory metrics;
%     anchor for our View B.
%   - SWE-bench family + integrity: SWE-bench / Verified / PolyBench / Pro /
%     SWE-bench+ / Illusion. The rising leak-audit bar.
%   - One sentence on SWE-Agent exclusion ("role in benchmark construction").
%   - One sentence on closed-source scope.
%
% Constraints:
%   - One paragraph per cluster, not a literature dump.
%   - Cite TRAJEVAL + SWE-Explore as the methodological anchor for View B.
% =========================================================================

\section{Related Work}
\label{sec:related}

\textbf{Coding-agent harnesses.} SWE-agent \citep{sweagent2024} introduced the
agent-computer-interface framing on top of a single-LLM control loop; OpenHands
\citep{openhands2025}, formerly OpenDevin, generalized the platform with sandboxed
execution and multi-agent coordination; Aider \citep{aider} drives file-level edits over a
local git repository with a dependency-ranked repo map; AutoCodeRover
\citep{autocoderover2024} pairs LLM reasoning with AST-aware code search and
spectrum-based fault localization; OpenCode \citep{opencode2025} is the model-agnostic
open-source TUI agent we use as the cross-harness comparator
(\S\ref{sec:system-opencode}). SuperCoder, the harness this paper studies
(\S\ref{sec:system-harness}), shares the parallel-tool-dispatch loop posture with these
systems but ships a structural codebase index as a first-class tool, which prior
measured harnesses do not. We exclude SWE-agent from the comparator set because its
agent-computer-interface pipeline was used in SWE-bench's construction, creating
circularity for an evaluation on SWE-bench-family tasks. Closed-source harnesses (Claude
Code, Cursor, Windsurf) are out of scope by design; this study fixes the model
(\S\ref{sec:design-model}) and varies the open-source harness configuration around it.

\textbf{Retrieval approaches for code agents.} Four approach types appear in the recent
literature. \emph{Agentic grep and read} drives OpenCode \citep{opencode2025} and similar
terminal-loop agents that call ripgrep over the working copy; there is no structural
codebase index. Sen et al.\ \citep{grepallyouneed2026} contrast grep with vector retrieval
inside agentic loops on the LongMemEval memory-retrieval benchmark, with grep favored;
their setting is non-code and the alternative they ablate against is dense retrieval, not
a structural codebase index, but the question framing (does grep suffice inside an
agentic harness?) is the closest prior to ours. \emph{Repository-level retrieval and
planning} approaches predate the agent-harness wave: RepoCoder \citep{repocoder2023}
iteratively retrieves over the whole repository for code completion, and CodePlan
\citep{codeplan2024} stages multi-file edits as a planned sequence of repository-wide
operations --- neither is built as an agent loop. \emph{File-dependency repo maps},
exemplified by Aider's PageRank-ranked repo map
\citep{aider}, surface candidate files by import and reference structure but do not index
symbol-level semantics. \emph{Semantic and graph search} combines code-chunk embeddings
with typed repository graphs: LocAgent \citep{locagent2025} equips an LLM agent with
graph-search tools over a heterogeneous code graph, RepoGraph \citep{repograph2024} plugs
a repository-wide code graph into SWE-agent and AutoCodeRover, and the Code Graph Model
line \citep{cgm2025} integrates the graph directly into an LLM's attention via an
adapter; Agentless \citep{agentless2024} reaches comparable SWE-bench-Lite scores with a
non-agentic three-phase localization-and-repair pipeline. \emph{Structural codebase
indices} have been adopted in commercial coding-agent stacks; this paper provides the
first leak-audited, model-controlled causal ablation of one inside an open-source harness
(\S\ref{sec:results-ablation}).

\textbf{Localization metrics for code agents.} LocAgent \citep{locagent2025} and
Agentless \citep{agentless2024} report file-level acc@$k$ as the primary localization
metric. The field has been moving toward stage-decomposed trajectory metrics: TRAJEVAL
\citep{trajeval2026} decomposes agent trajectories into search, read, and edit phases
with per-stage precision and recall, and SWE-Explore \citep{sweexplore2026} isolates
repository exploration as a sub-task with coverage and ranking metrics against
trajectory-derived ground truth. Our View~B
(\S\ref{sec:design-localization}, \S\ref{sec:results-crossharness}) sits in the same
field move: we strip engine-result paths from the surfaced set so that an SC-ON acc@$k$
counts the same kind of agent-targeted surface as an OpenCode acc@$k$. This paper does
not propose a new localization metric; it adopts the field-trend rule and applies it
uniformly across arms.

\textbf{SWE-bench family and benchmark integrity.} SWE-bench \citep{swebench2024}
introduced the canonical 2{,}294-issue Python benchmark; SWE-bench Verified
\citep{swebenchverified2024} released a 500-issue human-screened subset. Two recent
benchmark-expansion lines extend coverage: SWE-PolyBench \citep{swepolybench2025} is a
multi-language SWE-bench-style benchmark with a Verified subset, and SWE-bench Pro
\citep{swebenchpro2025} provides a harder long-horizon set; we use the Verified subset of
SWE-PolyBench and the public subset of SWE-bench Pro
(\S\ref{sec:design-benchmarks}). Three integrity audits motivate our hardening protocol:
SWE-bench+ \citep{swebenchplus2024} documented solution leakage in issue text and
weak-test pass-throughs in the original SWE-bench; the SWE-bench Illusion paper
\citep{swebenchillusion2025} showed that strong scores partly reflect memorization of
in-benchmark repositories; and Garg et al.\ \citep{savingswebench2025} mutate the formal
GitHub-issue specification into realistic user-style queries derived from chat-agent
telemetry, and report capability overestimation above 50\% for some models on
public benchmarks --- a phrasing-ecological-validity threat that detect-and-exclude
hardening (this paper) does not address. The per-cell scrub, fail-closed gate, and S1
reviewer pass in \S\ref{sec:integrity} sit in this audit tradition; the public exclusion
ledger (\S\ref{sec:integrity-ledger}) extends it by releasing per-cell evidence for every
dropped cell.
